%Dokumentenklasse
\documentclass[
	a4paper,										%Format
	12pt,											%Schriftgröße
	twoside,										%Zweiseitig
	headsepline,								%horizontale Linie oben
	footsepline,									%horizontale Linie unten
	captions=tableheading,					%Überschrift über Tabelle (nicht unter)
]{scrartcl}

\renewcommand{\thesubsection}{\thesection.\arabic{subsection}.}
\renewcommand{\thesubsubsection}{\thesection.\arabic{subsection}.\arabic{subsubsection}. \hspace{2pt}}
\renewcommand{\theparagraph}{\thesection.\arabic{subsection}.\arabic{subsubsection}.\arabic{paragraph}.}

%Pakete:
\usepackage[exponent-product = \cdot, output-decimal-marker = {,}]{siunitx}
\usepackage{textalpha}
\usepackage{floatrow}
\floatsetup[table]{style=plaintop}
\floatsetup[longtable]{LTcapwidth=table}
%\usepackage{overcite}
\usepackage{siunitx}
\usepackage[ngerman]{babel} 			%Neue deutsche Rechtschreibung
\usepackage{grffile}
\usepackage{graphicx} 					%Bilder einbinden
\usepackage{xcolor}						%Farben für Schrift und Plots
\usepackage{pgfplots}
\pgfplotsset{compat=1.18}					%Plots und Tabellen aus Datensätzen.
\usepackage{amsmath}						%Mathe
\usepackage{amssymb}					%Mathesymbole
\usepackage{longtable}					%Seitenübergreifende Tabellen zulassen
\usepackage{float}
\usepackage{tabularx}
\usepackage{geometry}					%Seitenränder
\usepackage{ae}								%skalierbare Schriften (im PDF-File)
\usepackage[activate]{microtype}		%optischer Randausgleich
\usepackage{setspace}					%Zeilenabstand: \onehalfspacing
\usepackage[utf8]{inputenc}				%Zeichenkodierung für Umlaute
\usepackage{scrlayer-scrpage}					%für Dokumentformatierung mit Pagestyle
\usepackage[version=4]{mhchem}						%Chemische Formeln (mit \ce )
\usepackage{chemfig}
\usepackage{siunitx}
\usepackage{caption}										%Überschriften formatieren
\captionsetup{labelfont={bf,it},tablename=Tab.,figurename=Abb.}				%Überschriften formatieren
\captionsetup[table]{singlelinecheck=false,justification=raggedright}
\captionsetup[figure]{justification=raggedright}				

%Abbildungs und Tabellennummerierung anpassen
\renewcommand*\figureformat {\figurename~\arabic{section}.\arabic{subsection}.\thefigure.}	%Punkt nach Abb. #
\renewcommand*\tableformat {\tablename~\arabic{section}.\arabic{subsection}.\thetable:}		%Doppelpunkt nach Tab. #
\renewcommand*\captionformat {~}								%Nur Abstand, kein Zeichen nach Tab. #: bzw Abb. #. 

%Literatur als Subsubsection definieren und von Literaturverzeichnis in Literatur umbenennen
\makeatletter
\renewcommand*\bib@heading{\subsubsection{\bibname}			
\@mkboth{\subsubsectionmarkformat \bibname}{\subsubsectionmarkformat \bibname}}
\makeatother
\addto\captionsngerman{\renewcommand*\bibname{Literatur}} 

\usepackage{prettyref}						%Nummerierung der Tabellen und Abbildungen im Text an Tab. x.x.# anpassen
%%% Für Abbildungen %%%
\newrefformat{fig}{Abb.~\arabic{section}.\arabic{subsection}.\ref{#1}}
%%% Für Tabellen %%%
\newrefformat{tab}{Tab.~\arabic{section}.\arabic{subsection}.\ref{#1}}

%Links im Dokument:
\usepackage[									%als letztes Paktet einbinden!
	linkcolor=black,							%Farbe der Links
	citecolor=black,							%Farbe der Link-Ziele
	urlcolor=black,								%Farbe von URLs
	colorlinks=true,							%Links nicht als Kasten, sondern direkt als Text anzeigen
]{hyperref}										%Links im PDF-File, unterstützung von URLs mit /url{}

%Dokumentformatierungen
\setkomafont{sectioning}{\bfseries} 	%Serifenschrift in Überschriften
\setkomafont{subsection}{\LARGE}
\setkomafont{subsubsection}{\large}
\pagenumbering{arabic}					%arabische Seitenzahlen
\geometry{left=3cm, right=2cm, top=2.5cm, bottom=2cm} 	%Seitenränder
\parindent0pt 									%Absätze einrücken um 0Pt
\parskip 10pt									%Abstand nach Absatz
\onehalfspacing								%Zeilenabstand
\pagestyle{scrheadings}					%Kopf- und Fussleistenformatierung
\headsep0.4cm								%Abstand Kopfleiste-Text
\footskip0.8cm									%Abstand Fussleiste-Text
\urlstyle{same}								%URLs normal formatieren, nicht als Typewriter

%Befehle definieren

%%%%% create chemequation environment

\newcounter{chemequation}
\renewcommand*\thechemequation{(\Roman{chemequation})}
\newenvironment{chemequation}{%
	\stepcounter{chemequation}%
	\begin{equation}}%
{\tag*{\thechemequation}%
	\end{equation}}

%%%%% create chemtag to iterate and ref chemequations in an align environment

\newcommand*\chemtag{%
	\stepcounter{chemequation}%
	\tag*{\thechemequation}}

%%%%% to create a chemicalequation float

\usepackage{float}
\newfloat{chemicalequation}{htbp}{cex}%[chapter]
\floatname{chemicalequation}{}

\newcommand{\zitat}[1]{\textsuperscript{\cite{#1}}}

\newcommand{\bild}[4]{					%Bilder mit Beschriftung einfügen
	\begin{figure}[#4]						%#4: Platzierung (htbp)
		\centering								%Bilder mittig einfügen
		\captionbox								%Beschriftungsumgebung
			{#2\label{fig: #1}}				%#2: Beschriftung
			{\includegraphics[#3]{#1}}	%#1: Dateiname, #3: Bildeigenschaften
	\end{figure}} 

\newcommand{\formel}[2]{				%Formel einfügen und Referenz erzeugen
	\begin{equation}
		\label{eq: #1}							%#1: Formelname
		#2											%Formelinhalt
	\end{equation}}

\newcommand{\cformel}[2][]{				%Chemische Formel einügen und Referenz erzeugen
	\begin{chemequation}
		\label{ceq: #1}					%#1: Formelname (optional)
		\ce{#2}							%#2: Formelinhalt
	\end{chemequation}\\}

\newcommand{\ox}[2]{
	\overset{\text{#2}}{\ce{#1}}}
	
\newcommand{\tabelle}[6]{				%Tabelle mit Beschriftung einfügen
\begin{table}[#3]							%#3: Platzierung (htbp)
\centering										%Tabelle mittig einfügen
\captionbox										%Beschriftungsumgebung
{#2												%#2: Tabellenbeschriftung
\label{tab: #1}}								%#1: Tabellenname
{\begin{tabular}{#4}						%#4: Anzahl und Ausrichtung Spalten (l c r)
\hline 											%horizontale Linie
#5													%#5: Obere Zeile
\hline
#6													%#6: Datenzeilen
\hline 
\end{tabular}}
\end{table}}

\usepackage{array,ragged2e}
\newcolumntype{P}[1]{>{\RaggedRight\arraybackslash}p{#1}} %Verhindert ungewollte Lücken in Tabellen

\newcommand{\langetabelle}[6]{  %Seitenübergreifende Tabelle, Kopfzeile wird 
	\begin{centering}           %auf jeder neuen Seite ebenfalls angezeigt
		{\begin{longtable}[#3]{#4}
				\caption{#2\label{tab: #1}}\\				
				\hline
				#5
				\hline
				\endfirsthead
                \caption{#2\label{tab: #1}}\\
				\hline
				#5
				\hline
				\endhead
				#6												
				\hline 
		\end{longtable}}
\end{centering}}

%Plots erstellen&&&&&&&&

%#1: Referenz
%#2: Name
%#3: x label
%#4: y label
%#5: x-werte,y-werte (Spaltenweise)
%#6: Quelle (.txt Datei)
%#7: "Nummerierung" (manuell) <.<

\newcommand{\plot}[6]{
	\begin{figure}[H]	
		\centering
		\begin{tikzpicture}
	
		\begin{axis}[
		title style={at={(0.5,0)},
		anchor=north,
		yshift=-40},
		xlabel={#3},
		ylabel={#4},
		legend pos=south east,
		legend entries={},
		]
		\addplot[color=black,mark=*,only marks] table [#5] {#6};
		\label{plot: plot1}
		\end{axis}
		\end{tikzpicture}
		\caption{#2}
		\label{fig: #1}
	\end{figure}
}

\newcommand{\molL}[1]{\ensuremath{\mathrm{\SI{#1}{\frac{mol}{L}}}}} % \molL{1} => \SI{1}{\frac{mol}{L}}
\newcommand{\mL}[1]{\ensuremath{\mathrm{\SI{#1}{mL}}}} % \mL{1} => \SI{1}{mL}

\newcommand{\Rn}[1]{\romannumeral#1\relax}                   %\Rn{1} = i
\newcommand{\RN}[1]{\uppercase\expandafter{\romannumeral#1}} %\RN{1} = I

%Dokument:%%%%%%%%%%%%%%%%%%%%%%%%%%%%%%%%%%%%%%%%%%%%%%%%%%%%%%%%%%%%%%%%%%%%%%%%%%%%
\begin{document}
	
	%Titelseite
	\begin{titlepage}
		\begin{center}
			\hrule													%waagerechte Linie
			\vspace*{5cm}
			\Huge {\textbf {Versuchsreihe 3}}\\						%Titel Versuchsreihe
			\vspace*{0.5cm}
			%		\large {\bf {Untertitel}}\\						%Untertitel
			\vspace*{1cm}
			\normalsize
			Protokoll	\\									
			Praktikum für Nanowissenschaftler:innen\\
			im Studiengang Nanowissenschaften\\
			\vspace*{0.5cm}
			\Large {\textbf{Michael Quint}}		\\						%Name
			\vspace*{0.5cm}
					\normalsize{Platz 9a}								%Platznummer
			\vfill
		\end{center}
		\begin{tabbing}
			\hspace{3cm}\=\kill
			Assistent: \> Dr. Hauke Heller \\						% Name Assistent
			Version: \> 1 \\									%Version
			Datum: \> 25.10.2026 \\								%Datum
		\end{tabbing}
		\hrule														%Waagerechte Linie
\end{titlepage}

\setcounter{page}{1}
\setcounter{section}{3}								%Versuchstag #.x.x bzw Versuchsreihe
\setcounter{subsection}{0}							%Versuchsnummer x.#.x startet bei 0 für 1 
\setcounter{secnumdepth}{4}
	
\subsection{Verschieben eines heterogenen Gleichgewichts}

\subsubsection{Einleitung und Aufgabenstellung}

	Das Verhalten einer gesättigten Kaliumperchloratlösung (\ce{KClO4}) sollte nach Zugabe einer gesättigten Natriumchlorid- (\ce{NaCl}),
	Kaliumchlorid- (\ce{KCl}) und 70\%iger Perchlorsäurelösung (\ce{HClO4}) untersucht und beschrieben werden. Die Begriffe eines homogenen und heterogenen Gleichgewichts waren zu definieren.
	Die Löslichkeit von Kaliumperchlorat war zu berechnen und für 100 mL Wasser anzugeben.

\subsubsection{Theoretischer Hintergrund}

	Chemische Gleichgewichte sind Zustände innerhalb chemischer Reaktionen in welchen keine Nettostoffmenge umgesetzt wird. Dies bedeutet nicht, dass keine Reaktion stattfindet,
	jedoch dass sowohl die Hin- als auch Rückreaktion die gleiche Reaktionsgeschwindigkeit vorweisen, sofern keine äußere Einwirkung vorliegt.\\
	Grundlegend können chemische Reaktionen mit dem Edukt \ce{A_i} und dem Produkt \ce{B_j} mit dem jeweiligen stöchiometrischen Faktor $\nu$ beschrieben werden.

	\cformel[Reaktion1]{\sum_{i} $\nu_i$ A_i <=> \sum_{j} $\nu_j$ B_j}
	{\centering
	$\nu_i > 0$, $\nu_j > 0$
	\par}
	Nach dem Prinzip von Le Chatelier ist dieses Gleichgewicht abhängig von äußeren Zwängen; dazu gehören die Temperatur, die Konzentrationen der Reaktanten und der Umgebungsdruck.
	Dabei weicht das Gleichgewicht jenem Zwang aus, dem es ausgesezt ist.\zitat{riedel}\\\\
	Es wird außerdem zwischen einem homogenen und heterogenen Gleichgewicht unterschieden. In einem homogenen Gleichgewicht befinden sich alle relevanten Stoff in derselben Phase,
	während in einem heterogenen Gleichgewicht mindestens zwei verschiedene Phasen vorliegen.\zitat{riedel}\\

	Mit dem Massenwirkungsgesetz (MWG) besitzen Gleichgewichtsreaktionen \prettyref{ceq: Reaktion1} eine Reaktionskonstante $K$, die folgende mathematische Beziehung aufweisen:

	\formel{MWG}{K = \frac{\prod_{j} a(\mathrm{B}_j)^{\nu_j}}{\prod_{i} a(\mathrm{A}_i)^{\nu_i}}}

	\newpage

	Dabei ist \ce{$a$(B_j)} die Aktivität des Produktes \ce{B_j} und \ce{$a$(A_i)} entsprechend die Aktivität des Eduktes \ce{A_i}.\\
	
	Bei der Dissoziation eines Salzes \ce{A_{x}B_y} in einer wässrigen Lösung gilt demnach folgendes Gleichgewicht:
	\cformel[Salzc]{{A_{x}B_y}_{(s)} <=> xA^{+}_{(aq)} + yB^{-}_{(aq)}}
	
	Folglich lässt sich $K$ wie folgt ausdrücken:
	\formel{Salzf}{K = \frac{a(\mathrm{A}^{+})^{x} \cdot a(\mathrm{B}^{-})^{y}}{a(\mathrm{A}_{x}\mathrm{B}_y)}}

	Da die Aktivität eines Feststoffes per Definition 1 beträgt und in dem Fall die Aktivität der Ionen näherungsweise gleich den Konzentrationen gesetzt wird, wird das Löslichkeitsprodukt $K_L$ eingeführt.
	\formel{KL}{K_L = [\mathrm{A}^{+}]^{x} \cdot [\mathrm{B}^{-}]^{y}}

	Für das Lösen eines Salzmoleküls entstehen $x+y$ Ionen \prettyref{ceq: Salzc}. Mit der Einführung der Löslichkeit $L$ lässt sich also ausdrücken:
	\formel{KL2}{K_L = (x\cdot L)^{x} \cdot (y\cdot L)^{y}}
	\formel{KL3}{L = \sqrt[x+y]{\frac{K_L}{x^x \cdot y^y}}}

	Angewendet auf Kaliumperchlorat bei einer Löslichkeitskonstante $K^{298\mathrm{K}}_L = 1,05 \cdot 10^{-2} \mathrm{\frac{mol^{2}}{L^{2}}}$ lässt sich die Löslichkeit $L^{298\mathrm{K}}_{\ce{KClO4}}$ berechnen.\zitat{KLWert}
	\cformel[KCL]{KClO4_{(s)} <=> K^{+}_{(aq)} + ClO4^{-}_{(aq)}}
	\formel{KL3}{L^{298\mathrm{K}}_{\ce{KClO4}} = \sqrt[1+1]{\frac{K^{298\mathrm{K}}_L}{1^1 \cdot 1^1}}}
	\formel{KL3}{L^{298\mathrm{K}}_{\ce{KClO4}} = \sqrt{K^{298\mathrm{K}}_L}}
	\formel{KL3}{L^{298\mathrm{K}}_{\ce{KClO4}} = \sqrt{1,05 \cdot 10^{-2} \mathrm{\frac{mol^{2}}{L^{2}}}} = 0,102\ \mathrm{\frac{mol}{L}} = 0,0102\ \mathrm{\frac{mol}{100\ mL}}}




\subsubsection{Experimenteller Teil}

\subsubsection{Ergebnisse und Diskussion}

\subsubsection{Zusammenfassung}

\newpage
\subsubsection{Sicherheit und Entsorgung}
		
Die für den Versuch verwendeten Chemikalien sind in \prettyref{tab: Chemikalien} mit Gefahrensymbolen, H- und P-Sätzen sowie Hinweisen zur Entsorgung aufgeführt.\zitat{test} Verwendete KMR-Stoffe sind in \prettyref{tab: KMR_leer} aufgeführt.\zitat{test}
	
\langetabelle
	{Chemikalien}
	{Verwendete Chemikalien.}
	{H}  % hier ggf. ein großes H eintragen, wenn die Tabelle sonst im Anhang verschwindet
	{l p{1.9cm} p{2.2cm} p{4.5cm} c}
	{Substanz & Symbol & H-Sätze & P-Sätze & Entsorgung \\}
	{ % Erster Stoff
  \ce{KClO4} 
    & \begin{tabular}[t]{l} GHS03 \\ GHS07 \end{tabular}
    & 271, 302 
    & 210, 220, 221, 264, \newline 280, 283, 301+312, \newline 306+360, 330, 370+378, \newline 371+380+375, 501 
    & (1) \\

  % Zweiter Stoff
  \ce{HClO4} (konz.) 
    & \begin{tabular}[t]{l} GHS03 \\ GHS05 \end{tabular}
    & 271, 314 
    & 210, 220, 221, 260, 264, \newline 280, 283, 301+330+331, \newline 303+361+353, 304+340, \newline 305+351+338, 306+360, \newline 310, 321, 370+378, \newline 371+380+375, 405, 501 
    & (1) \\}

\langetabelle
	{KMR_leer}
	{Verwendete kanzerogene, mutagene oder reproduktionstoxische Gefahrstoffe (KMR).}
	{H}
	{l P{3.8cm} P{2.4cm} P{3.5cm} c}
	{Substanz & H-Sätze & KMR-Einstufung & Verwendete Menge & Entsorgung \\}
	{ - & - & - & - & - \\}
	
Entsorgung: \\
\begin{itemize}
    \item[(1)] Alle verwendeten Lösungen wurden mit Natriumthiosulfat (\ce{Na2S2O3}) reduziert und im Behälter für nicht-nitrathaltige Säuren entsorgt.
\end{itemize}	

\newpage

\begin{thebibliography}{4}		%manuelles Quellenverzeichnis

	\bibitem{riedel}
						{A. Autor}, 
			\textit	{Titel}, 
						{Auflage}, 
						{Verlag}, 
						{Ort}, 
			\textbf	{Jahr}, 
			S.			{Seite-Seite}.

	\bibitem{KLWert}
						{Hans Lohninger}, 
						„{Tabelle von Löslichkeitsprodukten - Teil 2}“, zu finden unter 
			\url		{http://anorganik.chemie.vias.org/loeslichkeitsprodukt_tabelle_teil2.html}, 
			\textbf	{24.10.2026}.

	\bibitem{test} %Diese Quelle muss Jahrspezifisch angepasst werden
						{H. Heller}, 
			\textit	{Grundlagenpraktikum der Chemie für Studierende der Nanowissenschaften}, 
						{Version 1.22}, 
						{Universität Hamburg}, 
						{Hamburg}, 
			\textbf	{2026}, 
			S.			{34}.

			
\end{thebibliography}

%---------------------------------------------------------------------------------------------------------------------------------
\newpage

\setcounter{subsection}{1}							%Versuchsnummer x.#.x startet bei 0 für 1 
\setcounter{secnumdepth}{4}
	
\subsection{Fällung und Komplexierung von Silberionen}

\subsubsection{Einleitung und Aufgabenstellung}

\subsubsection{Theoretischer Hintergrund}

\subsubsection{Experimenteller Teil}

\subsubsection{Ergebnisse und Diskussion}

\subsubsection{Zusammenfassung}

\newpage
\subsubsection{Sicherheit und Entsorgung}
		
Die für den Versuch verwendeten Chemikalien sind in \prettyref{tab: Chemikalien2} mit Gefahrensymbolen, H- und P-Sätzen sowie Hinweisen zur Entsorgung aufgeführt.\zitat{hauke2} Verwendete KMR-Stoffe sind in \prettyref{tab: KMR_leer2} aufgeführt.\zitat{hauke2}
	
\langetabelle
	{Chemikalien2}
	{Verwendete Chemikalien.}
	{H}  % hier ggf. ein großes H eintragen, wenn die Tabelle sonst im Anhang verschwindet
	{l p{1.9cm} p{2.2cm} p{4.5cm} c}
	{Substanz & Symbol & H-Sätze & P-Sätze & Entsorgung \\}
	{ % Erster Stoff
  \ce{AgNO3} 
    & \begin{tabular}[t]{l} GHS03 \\ GHS05 \\ GHS09 \end{tabular}
    & 272, 314, \newline 410 
    & 210, 220, 221, 260, \newline 264, 273, 280, \newline 301+330+331, \newline 303+361+353, \newline 305+351+338, 310, \newline 321, 370+378, 391 \newline 405, 501, 234, 390 
    & (1) \\

  % Zweiter Stoff
  \ce{HCl} (0,1 mol/L) 
    & \begin{tabular}[t]{l} GHS05 \end{tabular}
    & 290
    & 234, 390
    & (1) \\
    
  \ce{NH4Cl} 
    & \begin{tabular}[t]{l} GHS07 \end{tabular}
    & 302, 319 
    & 264, 280, 301+312, \newline 305+351+338, 330, \newline 337+313, 501 
    & (1) \\

  \ce{BaCl2} 
    & \begin{tabular}[t]{l} GHS06 \end{tabular}
    & 301, 332 
    & 261, 264, 301+310, \newline 304+340, 321, 330, \newline 405, 501 
    & (1) \\

  \ce{KBr} 
    & \begin{tabular}[t]{l} GHS07 \end{tabular}
    & 315, 319, \newline 335 
    & 261, 264, 280, 302+352, \newline 304+340, 305+351+338, \newline 312, 321, 332+313, \newline 337+313, 362, 403+233, \newline 405, 501 
    & (1) \\
    
  \ce{KI} 
    & \begin{tabular}[t]{l} GHS08 \end{tabular}
    & 372 
    & 314 
    & (1) \\

  \ce{HNO3} (1 mol/L) 
    & \begin{tabular}[t]{l} GHS05 \end{tabular}
    & 314 
    & 260, 264, 280, \newline 301+330+331, \newline 303+361+353, 304+340, \newline 305+351+338, 310, 321, \newline 405, 501 
    & (1) \\

  \ce{C2H3Cl3O2} 
    & \begin{tabular}[t]{l} GHS06 \end{tabular}
    & 301, 315, \newline 319 
    & 264, 280, 301+310, \newline 302+352, 305+351+338, \newline 321, 330, 332+313, \newline 337+313, 362, 405, 501 
    & (1) \\

  \ce{KClO3} 
    & \begin{tabular}[t]{l} GHS03 \\ GHS07 \\ GHS09 \end{tabular}
    & 271, 302, \newline 332, 411 
    & 210, 220, 221, 261, 264, \newline 280, 283, 301+312, \newline 304+340, 306+360, 330, \newline 370+378, 371+380+375, \newline 391, 501 
    & (1) \\

  \ce{NaOH}
    & \begin{tabular}[t]{l} GHS05 \end{tabular}
    & 290, 314 
    & 260, 264, 280, \newline 301+330+331, \newline 303+361+353, 304+340, \newline 305+351+338, 310, 321, \newline 405, 501 
    & (1) \\

  \ce{AgCl} 
    & \begin{tabular}[t]{l} GHS05 \\ GHS09 \end{tabular}
    & 290, 400, \newline 410 
    & 234, 273, 390, 406, 501 
    & (1) \\

  \ce{AgBr} 
    & \begin{tabular}[t]{l} GHS09 \end{tabular}
    & 400, 410 
    & 273, 391, 501 
    & (1) \\

  \ce{AgI} 
    & \begin{tabular}[t]{l} GHS09 \end{tabular}
    & 400, 410 
    & 273, 391, 501 
    & (1) \\

  \ce{NH3} (25\,\%) 
    & \begin{tabular}[t]{l} GHS05 \\ GHS07 \\ GHS09 \end{tabular}
    & 290, 314, \newline 335, 400 
    & 260, 264, 273, 280, \newline 301+330+331, \newline 303+361+353, 304+340, \newline 305+351+338, 310, 321, \newline 391, 405, 501 
    & (1) \\

    }

\langetabelle
	{KMR_leer2}
	{Verwendete kanzerogene, mutagene oder reproduktionstoxische Gefahrstoffe (KMR).}
	{H}
	{l P{3.8cm} P{2.4cm} P{3.5cm} c}
	{Substanz & H-Sätze & KMR-Einstufung & Verwendete Menge & Entsorgung \\}
	{ - & - & - & - & - \\}
	
Entsorgung: \\
\begin{itemize}
    \item[(1)] 
    \item[(2)] 
\end{itemize}	

\newpage

\begin{thebibliography}{4}		%manuelles Quellenverzeichnis

	\bibitem{hauke2} %Diese Quelle muss Jahrspezifisch angepasst werden
						{H. Heller}, 
			\textit	{Grundlagenpraktikum der Chemie für Studierende der Nanowissenschaften}, 
						{Version 1.22}, 
						{Universität Hamburg}, 
						{Hamburg}, 
			\textbf	{2026}, 
			S.			{36-38}.

\end{thebibliography}

%--------------------------------------------------------------------------------------------------------------------------------------
\newpage

\setcounter{subsection}{2}							%Versuchsnummer x.#.x startet bei 0 für 1 
\setcounter{secnumdepth}{4}
	
\subsection{Komplexe des Kupfers}

\subsubsection{Einleitung und Aufgabenstellung}

\subsubsection{Theoretischer Hintergrund}

\subsubsection{Experimenteller Teil}

\subsubsection{Ergebnisse und Diskussion}

\subsubsection{Zusammenfassung}

\newpage
\subsubsection{Sicherheit und Entsorgung}
		
Die für den Versuch verwendeten Chemikalien sind in \prettyref{tab: Chemikalien3} mit Gefahrensymbolen, H- und P-Sätzen sowie Hinweisen zur Entsorgung aufgeführt.\zitat{hauke3} Verwendete KMR-Stoffe sind in \prettyref{tab: KMR_leer3} aufgeführt.\zitat{hauke3}
	
\langetabelle
	{Chemikalien3}
	{Verwendete Chemikalien.}
	{H}  % hier ggf. ein großes H eintragen, wenn die Tabelle sonst im Anhang verschwindet
	{l p{1.9cm} p{2.2cm} p{4.5cm} c}
	{Substanz & Symbol & H-Sätze & P-Sätze & Entsorgung \\}
	{ % Erster Stoff
  \ce{HCl} (konz.) 
    & \begin{tabular}[t]{l} GHS05 \\ GHS07 \end{tabular}
    & 290, 314, \newline 335
    & 260, 264, 280, \newline 301+330+331, \newline 303+361+353, \newline 304+340, \newline 305+351+338, \newline 310, 321, 403+233, \newline 405, 501
    & (1) \\

  \ce{NH3} (25\,\%) 
    & \begin{tabular}[t]{l} GHS05 \\ GHS07 \\ GHS09 \end{tabular}
    & 290, 314, \newline 335, 400 
    & 260, 264, 273, 280, \newline 301+330+331, \newline 303+361+353, 304+340, \newline 305+351+338, 310, 321, \newline 391, 405, 501 
    & (1) \\

  \ce{CuSO4}
    & \begin{tabular}[t]{l} GHS07 \\ GHS09 \end{tabular}
    & 302, 315, \newline 319, 335
    & 264, 273, 280, 301+312, \newline 302+352, 305+351+338, \newline 321, 330, 332+313, \newline 337+313, 362, 391, 501
    & (1) \\
    
  \ce{CuCO3.Cu(OH)2}
    & \begin{tabular}[t]{l} GHS07 \end{tabular}
    & 302, 315, \newline 319, 335
    & 261, 264, 280, 301+312, \newline 302+352, 304+340, \newline 305+351+338, 321, 330, \newline 332+313, 337+313, 362, \newline 403+233, 405, 501
    & (1) \\

    }

\langetabelle
	{KMR_leer3}
	{Verwendete kanzerogene, mutagene oder reproduktionstoxische Gefahrstoffe (KMR).}
	{H}
	{l P{3.8cm} P{2.4cm} P{3.5cm} c}
	{Substanz & H-Sätze & KMR-Einstufung & Verwendete Menge & Entsorgung \\}
	{ - & - & - & - & - \\}
	
Entsorgung: \\
\begin{itemize}
    \item[(1)] 
    \item[(2)] 
\end{itemize}	

\newpage

\begin{thebibliography}{4}		%manuelles Quellenverzeichnis

	\bibitem{hauke3} %Diese Quelle muss Jahrspezifisch angepasst werden
						{H. Heller}, 
			\textit	{Grundlagenpraktikum der Chemie für Studierende der Nanowissenschaften}, 
						{Version 1.22}, 
						{Universität Hamburg}, 
						{Hamburg}, 
			\textbf	{2026}, 
			S.			{39}.

\end{thebibliography}

\end{document}

%\newpage

%%%%%%%%%%%%%%%%%%%%%%%%%%%%%%%%%%%%%%%%%%%%%%%%%%%%%%%%%%%%%%%
Neuer Versuch:
\subsection{Versuchsname}

\subsubsection{Einleitung und Aufgabenstellung}
\subsubsection{Theoretischer Hintergrund}
\subsubsection{Experimenteller Teil}
\subsubsection{Ergebnisse und Diskussion}
\subsubsection{Zusammenfassung}
\subsubsection{Sicherheit und Entsorgung}	
Die für den Versuch verwendeten Chemikalien sind in \prettyref{tab: entsorgung1} mit Gefahrensymbolen, H- und P-Sätzen sowie Hinweisen zur Entsorgung aufgeführt.\zitat{Gefahren1} Verwendete KMR-Stoffe sind in \prettyref{tab: KMR1} aufgeführt.
	
\langetabelle
	{entsorgung1}
	{Verwendete Chemikalien.}
	{H}  % hier ggf. ein großes H eintragen, wenn die Tabelle sonst im Anhang verschwindet
	{l p{1.4cm} p{2.2cm} p{4.5cm} c}
	{Substanz & Symbol & H-Sätze & P-Sätze & Entsorgung \\}
	{\\}
    

\langetabelle
	{KMR1}
	{Verwendete kanzerogene, mutagene oder reproduktionstoxische Gefahrstoffe (KMR).}
	{H}
	{l P{3.8cm} P{2.4cm} P{3.5cm} c}
	{Substanz & H-Sätze & KMR-Einstufung & Verwendete Menge & Entsorgung \\}
	{ - & - & - & - & - \\}
Entsorgung: \\
\begin{itemize}
    \item[(1)] 
\end{itemize}	

\newpage

\begin{thebibliography}{4}		%manuelles Quellenverzeichnis
	\bibitem{Gefahren1}
						{H. Heller}, 
			\textit	{Grundlagenpraktikum der Chemie für Studierende der Nanowissenschaften}, 
						{Version 1.18},
                        {Universität Hamburg},
						{Hamburg}, 
			\textbf	{2023}, 
			S.{34}.
			
			
\end{thebibliography} 
\newpage

Nützliche Beispiele:

%Tabelle für Reagenzgläserinhalte
\tabelle
	{Tabellenname}												%Eindeutige Bezeichnung der Tabelle. \prettyref{tab: Tabellenname}
		{Tabellenüberschrift.}	%Inhalt der Reagenzgläser ... bis...
		{htb}	
		{l l l}					
		{Reagenzglas & Inhalt & Beschreibung}				
		{1 & 2~mL NaOH & klare, farblose Lsg.\\
		2 & & 
		}										


%Formel Löslichkeit KClO4 -> chemische Formel
\begin{align}
    \chemtag
    \ce{KClO4(aq) <=> K+(aq) + ClO4-(aq)}
\end{align}

%hier allgemeine Formel 
\formel
		{formel2}				%eindeutige Bezeichnung der Formel, Verweis: \ref{eq: formelname}
		{L_{\mathrm{A}_{a} B_{b}} = \sqrt[m+n]{\frac{K_{\mathrm{L}}}{a^{a} \cdot b^{b}}}  }


Bilder:
	\bild
		{bildname}					%Dateiname des Bildes, eindeutiger Name des Bildes, Verweis: \prettyref{fig: Bildname}
		{Bildbeschreibung.		%Text, der unter dem Bild erscheint
		\zitat{Quelle}}			%Name der Quelle im Literaturverzeichnis
		{Bildeigenschaften}		%angle, draft, scale, height, width
		{Platzierung}				%h: here, t: top, b: bottom, p: page; bestimmt, wo das Bild eingefügt wird. p: spezielle Seite für Abbildungen und Tabellen, Kombinationen möglich, Versuch der Platzierung gemäß Reihenfolge.
	
		
	Bildeigenschaften:
		angle - Bestimmt den Drehwinkel um einen Referenzpunkt für die Grafik. Positive Werte drehen die Grafik gegen den Uhrzeigersinn, negative mit ihm. 
		draft - Ist draft gesetzt, wird nur ein Platzhalter der Grafik angezeigt. Beschleunigt den Übersetzungsvorgang des Dokuments erheblich. 
		scale - Skaliert die Grafik anhand eines Skalierungsfaktors. 
		height - Skaliert die Grafik auf die angegebene Höhe. Die Angabe der Größe muss in einer LaTeX-spezifischen Einheit erfolgen. 
		width - Skaliert die Grafik auf die angegebene Breite. Die Angabe der Größe muss in einer LaTeX-spezifischen Einheit erfolgen. 	
	

Formeln:
	\formel
		{formelname}				%eindeutige Bezeichnung der Formel, Verweis: \ref{eq: formelname}
		{Formel}						%


Tabellen:
	\tabelle
		{Tabellenname}												%Eindeutige Bezeichnung der Tabelle. \prettyref{tab: Tabellenname}
		{Tabellenüberschrift.\zitat{Quelle}}						%Text, der über der Tabelle erscheint, Quelle der Daten
		{Platzierung}													%h: here, t: top, b: bottom, p: page; bestimmt, wo das Bild eingefügt wird. p: spezielle Seite für Abbildungen und Tabellen, Kombinationen möglich, Versuch der Platzierung gemäß Reihenfolge.		
		{Spaltenanzahl_und_Ausrichtung}						%{ l c r }: 3 Spalten, links mittig und rechts; { c c c c }: 4 Spalten, alle mittig
		{Titelzeile}														%Obere Zeile(n), danach folgt ein waagerechter Strich, Spalten durch & trennen, Neue Zeile durch \\
		{Datenzeilen}													%Zeilen unter dem waagerechten Strich, Spalten durch & trennen, Neue Zeile durch \\
		
	Zellen verbinden:	
		\multicolumn{Anzahl_der_verbunden_Spalten}{Ausrichtung (l/c/r}{Inhalt}
	
Reaktionen:
Formeln: 
\begin{center}
\ce{NH+4}, \ce{Fe3(CO)12}, \ce{Al^{3+}},\ce{Cr2O^{2-}7}
\end{center}

Isotope:
\begin{center}
\ce{^{37}_{17}Cl}, \ce{^{32}_{16}S^{2-}}
\end{center}

Bindungen: 
\begin{center}
\ce{H-O-N=S}, \ce{Cl-C#N}
\end{center}

Reaktionsgleichungen: 
\begin{center}
\ce{2H2 + O2 -> 2H2O}
\ce{MgCl2 ->[\text{Elektrolyse}] Mg\,(l) + Cl_{2}\,(g)}
\end{center}

Redoxreaktion:
\begin{center}
$\overset{0}{\ce{S}} + \overset{0}{\ce{O2}} \ce{->} \overset{+\text{IV}}{\ce{S}}+\overset{-\text{II}}{\ce{O2}}$
\end{center}


Literaturverweise:
		\bibitem{zeitschrift}
						{A. Autor}, 
			\textit	{Zeitschrift} 
			\textbf	{Jahr}, 
			\textit	{Band} 
						{Seite-Seite}.
		
		\bibitem{buch_ohne_herausgeber}
						{A. Autor}, 
			\textit	{Titel}, 
						{Verlag}, 
						{Ort}, 
			\textbf	{Jahr}, 
			S.			{Seite-Seite}.
			
		\bibitem{buch_mit_herausgeber}
						{A. Autor}, 
			\textit	{Titel}, 
						{Auflage}, 
						{Verlag}, 
						{Ort}, 
			\textbf	{Jahr}, 
			S.			{Seite-Seite}.
			
		\bibitem{dissertation}
						{A. Autor}, 
						{Dissertation}, 
						{Universität}, 
			\textbf	{Jahr}.
			
		\bibitem{internetseite}
						{A. Autor}, 
						„{Titel}“, zu finden unter 
			\url		{www.xyz.de}, 
			\textbf	{Jahr}.

% Kommentare Sophie:

% <- Für Kommentare oder Code der gerade nicht gebraucht wird, wird beim kompilieren nicht beachtet
%Um Bilder einzufügen: \bild{'Dateiname' (Datei muss im gleichem Ordner mit der TeX Datei liegen)}{'Bildunterschrift'}{width= ...cm (Größe des Bildes)}{H (Ort des Bildes, H=Hier)}

%Für Summenformeln im Text: \ce{z.B.: H2SO4 (Zahlen werden automtisch tief gestellt)}

% Um Formeln im Text zu referieren: \prettyref{fig (für Bilder und Graphen)/ eq (für mathematische Formeln)/ ceq (für chemische Formeln): 'Referenz'}

%Für chemische Formeln: \cformel['Referenz']{2\ox{Fe}{+\RN{3}}Cl3 + \ox{Fe}{+\RN{2}}Cl2 + 8NH4+ + 8 OH- \ce{<=>[{N2H4}]} \ox{Fe}{+\RN{2}}O\cdot \ox{Fe}{+\RN{3}}2O3 \ce{v} + 4 H2O + 8NH4+ + 4 Cl-}

%Für mathematische Formeln: \formel{'Referenz'}{2\cdot d\cdot \ce{sin}(\theta)= n \lambda}

%Tabellen: Online Converter um Exel Tabellen in LaTeX Format zu bringen: https://www.tablesgenerator.com
%\tabelle{'Referenz'}{'Tabellenüberschrift'}{H}{'Breite der Spalte bzw. Textausrichtung' (Spaltenbreite passt sich Textlänge an: l=links, r=rechts, c=zentriert. Spaltenbreite fest und Blocksatz: p{...cm}), Jede Spalte muss einzeln definiert werden, z.B l l l = 3 Spalten}{'Spaltenüberschrift' muss der Anzahl an Spalten entsprechen. Spaltenüberschriften müssen durch ein & getrennt werden. Außerdem wichtig am Ende: \\}{'Tabelleneinträge', müssen ebenfalls durch ein & getrennt werden und am Ende der Zeile ein \\ haben
%	- & - & - & \\
%	- & - & - & \\}
%Beispiel für eine Tabelle (zum Ausführen '%' entfernen):
%\tabelle{Auflösungsvermögen}{Experimentell ermitteltes und theoretisches Auflösungsvermögen für Objektiv 20X und 63X bei schräger und gerader Beleuchtung.}{H}{l l p{2cm} p{3cm}  p{3cm}}{Objektiv&Beleuchtung& Spalt mit einem Peak&experimentelles Auflösungsvermögen~/~nm&theoretisches Auflösungsvermögen~/~nm\\}
%	{20X&gerade&16&>750&1125\\
%	20X&schräg&9&690&563\\
%	63X&gerade&11&>750&600\\
%	63X&schräg&7&584&300\\}